% ---------------------------------------------------------------------
%  06-practice.tex -- HPB in practice: canonical situations
%  Every condition figure is \input exactly once, next to its narrative.
%  Each situation follows the same pattern:
%    situation -> how HPB draws it -> decision it supports.
% ---------------------------------------------------------------------
\section{HPB in Practice: Canonical Situations}
\label{sec:practice}

A notation only becomes a shared language when the situations that keep
coming back are always drawn the same way. Running a portfolio of
hardware, firmware and software products, I found that a small number of
situations account for almost every planning and traceability question I
had to answer. This section walks through nine of them. Together they
exercise every rule R1--R7.

Each situation is told in three short steps: \emph{the situation} as a
product manager meets it, \emph{how HPB draws it}, and \emph{the decision}
the diagram supports. All products, items and dates are fictitious and
illustrative, and each figure stands on its own. Time runs left to right
in \emph{events}, not days (\cref{sec:notation-time}).
\Cref{tab:conditions} is the map of the section.

Read every diagram in the same order: the main line from left to right,
then above it (how the product gets better), then below it (what is wrong
and how it gets fixed), then the NOW marker, right of which everything is
plan. With several lines or streams, do this once per band. Then ask of
any branch the three questions HPB exists to answer: \emph{Where did it
start? What else was open at the same time? Which version closed it, or
why was it dropped?}

\begin{table}[t]
  \centering
  \small
  \caption{The nine canonical situations and the decision each diagram
    supports. All examples are illustrative.}
  \label{tab:conditions}
  \renewcommand{\arraystretch}{1.1}
  \begin{tabularx}{\linewidth}{@{}>{\raggedright\arraybackslash}p{3.6cm}
      >{\centering\arraybackslash}p{0.7cm}
      >{\raggedright\arraybackslash}X
      >{\raggedright\arraybackslash}p{1.2cm}@{}}
    \toprule
    \textbf{Situation} & \textbf{Fig.} & \textbf{Decision supported}
      & \textbf{Rules} \\
    \midrule
    A. Firmware or software line
      & \ref{fig:software-line}
      & Which version contains a change; in which version a defect was
        seen; release notes per version.
      & R1--R3 \\
    B. Release gate
      & \ref{fig:release-gate}
      & Can the planned version ship? Release, slip or descope.
      & R4 \\
    C. Hardware revision, new reference
      & \ref{fig:hardware-fork}
      & What goes into the new reference; how long the old one is
        maintained; which fixes must be carried over.
      & R5, R3 \\
    D. Discard and re-target
      & \ref{fig:discard}
      & Dropping work explicitly; moving a commitment without losing its
        history.
      & R2, R7 \\
    E. Many branches at once
      & \ref{fig:scale}
      & Which long-lived open work needs a decision; how much is in
        flight per type.
      & R6 \\
    F. Kanban board versus HPB
      & \ref{fig:same-data}
      & What a team board can and cannot tell you about the product.
      & R1, R2 \\
    G. Maintained versions, back-port
      & \ref{fig:backport}
      & Which maintained versions contain a fix; which fixes still have to
        be ported.
      & R2, R3 \\
    H. Field return (RMA)
      & \ref{fig:rma}
      & How a field failure was solved without a new design; whether the
        action worked; supplier-driven change.
      & R2, R7 \\
    I. Shared firmware, joint launch
      & \ref{fig:shared}
      & Which configurations a defect affects; whether a coordinated
        launch can go ahead.
      & R4, R1 \\
    \bottomrule
  \end{tabularx}
\end{table}

% =====================================================================
\subsection{A. A firmware or software line}
\label{sec:cond-software}

% F2 -- Condition A: a firmware/software product line (history only).
% Illustrative data. Requires hpb.sty v1.1. Canvas: document defaults.
\begin{figure}[t]
\centering
\hpbfit{%
\begin{hpbgraph}
  % lanes: improvements above, issues and tasks below the main line
  \hpblane{1}{imp}{Improvement (I)}{0}{16.7}
  \hpblane{-1}{iss}{Issue (X)}{0}{16.7}
  \hpblane{-2}{task}{Task (T)}{0}{16.7}
  % main line and baselines
  \hpbmainline{0}{16.7}
  \hpborigin{0}{Product A}
  \hpbrelease{1}{v1.0}
  \hpbrelease{5}{v1.0.1}
  \hpbrelease{10}{v1.1}
  \hpbrelease{16}{v1.2}
  % branches (fork col, lane, start col, end col, end, merge col)
  \hpbbranch[label={X-101 Boot loop}]{iss}{1}{-1}{2}{4}{merge}{5}
  \hpbbranch[forkoff=0.6, mergeoff=0.6, label={I-21 Energy-saving mode}]{imp}{1}{1}{3}{9}{merge}{10}
  \hpbbranch[forkoff=0.6, label={X-102 Log time zone}]{iss}{5}{-1}{6}{8}{merge}{10}
  \hpbbranch[forkoff=0.3, mergeoff=0.3, label={T-8 Renew signing certificate}]{task}{5}{-2}{7}{14}{merge}{16}
  \hpbbranch[forkoff=0.6, mergeoff=0.6, label={I-22 Remote diagnostics}]{imp}{10}{1}{11}{15}{merge}{16}
  \hpbbranch[mergeoff=0.6, label={X-103 Wi-Fi reconnect}]{iss}{10}{-1}{12}{13}{merge}{16}
  % ordinal time axis with calendar ticks (month of each baseline)
  \hpbaxis{0}{16.7}{-2.75}
  \hpbtick{1}{Jan 2026}{-2.75}
  \hpbtick{5}{Feb}{-2.75}
  \hpbtick{10}{Apr}{-2.75}
  \hpbtick{16}{Jul}{-2.75}
\end{hpbgraph}}
\caption{Situation~A: a firmware or software line, history
only (illustrative). Each branch shows the version it started from and the
version that closed it; X-101 is a hotfix closed by \texttt{v1.0.1}. Time is
ordinal: one column per event, with months only as ticks.}
\label{fig:software-line}
\end{figure}


\paragraph{Situation.}
One line, released in successive versions, with improvements, defects and
tasks running in parallel. A customer reports a boot loop: is it fixed, and
since when? Sales asks what went into \texttt{v1.1}. A backlog knows the
items, not the versions around them.

\paragraph{In HPB.}
\Cref{fig:software-line} is pure history: four releases of
\emph{Product~A}, no NOW marker. I-21 forks from \texttt{v1.0} and merges
into \texttt{v1.1}; its path passes over the patch \texttt{v1.0.1}, so the
patch visibly does not contain it. X-101 \emph{Boot loop} is seen in
\texttt{v1.0} and closed by the hotfix \texttt{v1.0.1}. T-8 stays open
across \texttt{v1.1} and closes with \texttt{v1.2}.

\paragraph{Decision.}
Every history question is one path. The boot loop is fixed in
\texttt{v1.0.1} and in everything to its right. Release notes are the
rises into each version; escaped defects are the drops below it. R1--R3
can be checked by eye.

% =====================================================================
\subsection{B. Planning with a release gate}
\label{sec:cond-gate}

% F3 -- Condition B: planning with a release gate.
% Illustrative data. Requires hpb.sty v1.1. Canvas: document defaults.
% Lanes follow Algorithm 1: closed branches take the innermost free lane,
% branches open at NOW the outermost one.
\begin{figure}[t]
\centering
\hpbfit{%
\begin{hpbgraph}
  % lanes: improvements above, issues below (three issue lanes: lane packing)
  \hpblane{1}{imp}{Improvement (I)}{0}{17.7}
  \hpblane{-1}{iss}{Issue (X)}{0}{17.7}
  \hpblane{-2}{iss}{Issue (X)}{0}{17.7}
  \hpblane{-3}{iss}{Issue (X)}{0}{17.7}
  % main line: history (solid) up to NOW, plan (dashed) afterwards
  \hpbmainline{0}{11}
  \hpbmainline*{11}{17.7}
  \hpborigin{0}{Product A}
  \hpbrelease{1}{v2.0}
  \hpbrelease{6}{v2.0.1}
  \hpbplanned{13}{v2.1}{May 2027}
  \hpbplanned{17}{v2.2}{Sep 2027}
  % history: a hotfix already merged (closed: innermost lane)
  \hpbbranch[forkoff=0.6, label={X-210 Sync error}]{iss}{1}{-1}{2}{5}{merge}{6}
  % open but untargeted issue (open at NOW: outermost lane); no gate
  \hpbbranch[forkoff=0.3, label={X-213 Display flicker (untargeted)}]{iss}{1}{-3}{4}{11}{open}{}
  % targeted branches: open at NOW, dashed continuation into the gate of v2.1
  % (nested rises: farther lane, smaller offset; offsets 0.2 column apart)
  \hpbbranch[forkoff=0.5, mergeoff=0.8, label={I-50 Multi-user access}]{imp}{1}{1}{3}{11}{target}{13}
  \hpbbranch[forkoff=0.3, mergeoff=0.4, label={X-211 Report totals}]{iss}{6}{-2}{7}{11}{target}{13}
  \hpbbranch[forkoff=0.6, mergeoff=0.6, label={X-212 CSV encoding}]{iss}{6}{-1}{8}{11}{target}{13}
  % steps: a done sub-task on I-50; X-212 verified in release candidate RC1
  \hpbstep{imp}{10}{1}{done}
  \hpbstep{iss}{9}{-1}{pass}
  % gate around the planned baseline v2.1 and NOW marker
  \hpbgate{13}{1}{-2}{gate v2.1}
  \hpbnow{11}{1}{-3}
  \hpbcallout{6.3,2.35}{v2.1 cannot be published until I-50 and X-211 are verified\\
    (as X-212 is), re-targeted or discarded (rule R4)}{12.6,1.62}
  % ordinal time axis with calendar ticks (history part)
  \hpbaxis{0}{17.7}{-3.75}
  \hpbtick{1}{Jan 2027}{-3.75}
  \hpbtick{6}{Feb}{-3.75}
  \hpbtick{11}{Apr}{-3.75}
\end{hpbgraph}}
\caption{Situation~B: a release gate (illustrative). Three
branches are committed to \texttt{v2.1}, and their dashed lines enter its
gate; only X-212 has a passed verification. By \cref{rule:gate},
\texttt{v2.1} cannot ship until I-50 and X-211 are verified, moved or
discarded. X-213 is open but committed to nothing.}
\label{fig:release-gate}
\end{figure}


\paragraph{Situation.}
Three things are promised for a release; two are not finished. Too often
this is discovered in release week, and the choice to ship, slip or cut is
made in a corridor and never written down.

\paragraph{In HPB.}
In \cref{fig:release-gate}, three open branches are targeted at the
planned \texttt{v2.1}: each stops at NOW and continues as a dashed line
into the shaded gate. Only X-212 carries a passed test on a release
candidate. X-213 is open with no target; its path just ends at NOW.

\paragraph{Decision.}
The diagram \emph{is} the agenda of the gate review. By \cref{rule:gate},
\texttt{v2.1} cannot ship while I-50 and X-211 are unverified; each one is
verified, moved to \texttt{v2.2} or discarded with a reason, and the
outcome (release, slip or descope) is recorded. X-213 goes to triage
before it becomes an eternal branch.

% =====================================================================
\subsection{C. A hardware revision that creates a new reference}
\label{sec:cond-hardware}

% F4 -- Condition C: hardware revision and line fork.
% Illustrative data. Requires hpb.sty v1.1. Canvas: document defaults.
% The hardware baseline Rev B is the FIRST node of the child line (drawn
% above its parent, in its own lane band); the lineage connector leaves the
% parent main line after v1.3, the baseline Rev B was derived from.
\begin{figure}[t]
\centering
\hpbfit{%
\begin{hpbgraph}
  % ---- child line DMS-4000 Rev B (y = 3.5), its own band and lane
  \hpblane{4.5}{imp}{Improvement (I)}{0}{19.7}
  \hpblane{3.5}{main}{DMS-4000\\Rev B}{0}{9}
  \hpbmainline{9}{19.7}[3.5]
  \hpbseparator{2.85}{-2.4}{19.7}
  % ---- parent line DMS-4000 Rev A (y = 0) and its lanes
  \hpblane{2}{lab}{Lab (L)}{0}{19.7}
  \hpblane{1}{imp}{Improvement (I)}{0}{19.7}
  \hpblane{-1}{iss}{Issue (X)}{0}{19.7}
  \hpbmainline{0}{19}
  \hpborigin{0}{DMS-4000\\Rev A}
  \hpbrelease{1}{v1.2}
  \hpbrelease{5}{v1.3}
  \hpbmilestone{10}{end of sale}
  \hpbrelease{11}{v1.4}
  \hpbrelease{16}{v1.4.1}
  \hpbeos{19}
  % ---- line fork: lineage connector from the parent main line (after v1.3)
  %      to the hardware baseline Rev B, first baseline of the child line
  \hpblinefork{9}{0}{3.5}
  \hpbhardware{9}{Rev B}[3.5]
  \hpbrelease{12}{v2.0}[3.5]
  \hpbrelease{18}{v2.1}[3.5]
  % lab evolution with test steps, and an improvement, both forked from v1.2
  % and merged into Rev B (cross-line merge; nested rises)
  \hpbbranch[forkoff=0.5, mergeoff=0.6, mergeline=3.5,
    label={L-5 Sensor board v2}]{lab}{1}{2}{2}{7}{merge}{9}
  \hpbbranch[forkoff=0.8, mergeoff=0.4, mergeline=3.5,
    label={I-30 Low-power mode}]{imp}{1}{1}{3}{8}{merge}{9}
  \hpbstep{lab}{4}{2}{fail}
  \hpbstep{lab}{6}{2}{pass}
  % work on the child line
  \hpbbranch[line=3.5, label={I-40 Remote unlock}]{imp}{12}{4.5}{13}{17}{merge}{18}
  % Rev A in maintenance: a defect observed in v1.4, fixed by firmware v1.4.1
  \hpbbranch[label={X-118 Relay chatter}]{iss}{11}{-1}{12}{14}{merge}{16}
  % ---- annotation
  \hpbcallout{12.6,1.6}{line fork (R5): parent(Rev B) = Rev A,\\
    Rev B derived from v1.3}{9,1.6}
  % ---- ordinal time axis with calendar ticks
  \hpbaxis{0}{19.7}{-1.85}
  \hpbtick{1}{Jan 2026}{-1.85}
  \hpbtick{5}{Mar}{-1.85}
  \hpbtick{9}{Jun}{-1.85}
  \hpbtick{16}{Sep}{-1.85}
  \hpbtick{19}{Dec}{-1.85}
\end{hpbgraph}}
\caption{Situation~C: a hardware revision that creates a new
reference (illustrative). L-5 and I-30 merge into the diamond Rev~B, which
starts the child line \emph{DMS-4000 Rev~B}; the lineage connector leaves the
parent after \texttt{v1.3}, the version Rev~B derives from
(\cref{rule:lineage}). The parent keeps its own life: end of sale, the fix
for X-118 in \texttt{v1.4.1}, and end of service (EoS).}
\label{fig:hardware-fork}
\end{figure}


\paragraph{Situation.}
A hardware change fails the interchangeability test, so manufacturing
needs a new reference. From then on there are two physical products in
the field and every question doubles. This is where the methods of
\cref{sec:gaps} leave the widest gap.

\paragraph{In HPB.}
In \cref{fig:hardware-fork}, the lab branch L-5 (one failed and one passed
test) and the improvement I-30 merge into the hardware diamond
\emph{Rev~B}. Rev~B starts a child line in its own band, linked to the
parent after \texttt{v1.3}. New work lives on the child; the parent passes
end of sale, ships the fix for X-118 in \texttt{v1.4.1} and ends with the
\textsf{EoS} bar.

\paragraph{Decision.}
What the new reference contains (L-5 and I-30, backed by the passed test),
how long the old one lives (its milestones and the work still forking
from it), and which fixes must be carried over: the fix for X-118 is
\emph{not} in Rev~B, so its status there must be recorded, and ported if
needed (situation~G).

% =====================================================================
\subsection{D. Discarded work and re-targeting}
\label{sec:cond-discard}

% F5 -- Condition D: discarded and not-viable work, re-targeting.
% Illustrative data. Requires hpb.sty v1.1. Canvas: document defaults.
% Lanes follow Algorithm 1 (closed: innermost free lane; open at NOW:
% outermost free lane).
\begin{figure}[t]
\centering
\hpbfit{%
\begin{hpbgraph}
  % lanes: improvements above; two issue lanes below (parallel defects)
  \hpblane{1}{imp}{Improvement (I)}{0}{18.7}
  \hpblane{-1}{iss}{Issue (X)}{0}{18.7}
  \hpblane{-2}{iss}{Issue (X)}{0}{18.7}
  % main line: history up to NOW (col 13), plan afterwards
  \hpbmainline{0}{13}
  \hpbmainline*{13}{18.7}
  \hpborigin{0}{Product A}
  \hpbrelease{1}{v3.0}
  \hpbrelease{8}{v3.0.1}
  \hpbplanned{15}{v3.1}{Mar 2027}
  \hpbplanned{18}{v3.2}{Jun 2027}
  % improvement judged not viable: discarded, never rejoins the line
  \hpbbranch[label={I-60 Solar panel}]{imp}{1}{1}{2}{5}{discard}{}
  \hpblabel[right]{5}{1}{\itshape not viable\\\itshape (unit cost)}
  % defect fixed by the patch baseline v3.0.1
  \hpbbranch[forkoff=0.3, label={X-52 Display flicker}]{iss}{1}{-2}{3}{7}{merge}{8}
  % same defect reported again: closed as duplicate
  \hpbbranch[forkoff=0.6, label={X-57 Screen flickers}]{iss}{1}{-1}{4}{5}{discard}{}
  \hpblabel[right]{5}{-1}{\itshape duplicate\\\itshape of X-52}
  % improvement targeted at v3.1
  \hpbbranch[label={I-61 Offline mode}]{imp}{8}{1}{9}{13}{target}{15}
  % defect re-targeted from v3.1 to v3.2: the decision is a recorded (done)
  % step on the branch with a label; former target faded, current target at
  % full strength
  \hpbbranch[label={X-63 OTA timeout}]{iss}{8}{-2}{10}{13}{target}{18}
  \hpbtarget[faded]{13}{-2}{15}
  \hpbstep{iss}{11}{-2}{done}
  \hpblabel[below center]{11}{-2}{re-targeted\\[-1pt]v3.1\,$\rightarrow$\,v3.2}
  % NOW marker
  \hpbnow{13}{1}{-2}
  % annotation of the former target
  \hpbcallout{16.0,-1.0}{former target\\(faded)}{14.5,-1.0}
  % ordinal time axis with calendar ticks
  \hpbaxis{0}{18.7}{-3.1}
  \hpbtick{1}{Jul 2026}{-3.1}
  \hpbtick{8}{Sep}{-3.1}
  \hpbtick{13}{Nov}{-3.1}
  \hpbtick{15}{Mar 2027}{-3.1}
  \hpbtick{18}{Jun}{-3.1}
\end{hpbgraph}}
\caption{Situation~D: discards and re-targeting (illustrative). I-60
and the duplicate X-57 end in gray crosses with their reasons and never
rejoin the line. X-63 was moved from \texttt{v3.1} to \texttt{v3.2}; the old
commitment stays visible as a faded dashed line.}
\label{fig:discard}
\end{figure}


\paragraph{Situation.}
A feature turns out too expensive, two reports describe the same defect,
an issue is pushed to the next release. In list tools all three end as
``closed'' or as an edited date, and six months later nobody knows why.

\paragraph{In HPB.}
In \cref{fig:discard}, I-60 is discarded: its path turns gray, never
returns to the line and carries the reason \emph{not viable (unit cost)}.
X-57 is discarded as \emph{duplicate of X-52}. X-63 was targeted at
\texttt{v3.1} and moved to \texttt{v3.2}; the old target remains as a
faded rise.

\paragraph{Decision.}
Decisions \emph{not} to deliver are records, not deletions
(\cref{rule:closure,rule:discussion}). Discards can be counted by reason,
and the faded rise shows which commitment moved, which is what an honest
gate-slip report needs (\cref{tab:metrics}).

% =====================================================================
\subsection{E. Many branches at once}
\label{sec:cond-scale}

% F6 -- Condition E: readability at scale (collapsed history, lane packing).
% Illustrative data. Requires hpb.sty v1.1. Canvas: document defaults.
% Lanes follow Algorithm 1 (closed: innermost free lane; open at NOW:
% outermost free lane).
\begin{figure}[t]
\centering
\hpbfit{%
\begin{hpbgraph}
  % two lanes per type; older open branches occupy the outer lanes
  \hpblane{4}{lab}{Lab (L)}{0}{19.7}
  \hpblane{3}{lab}{Lab (L)}{0}{19.7}
  \hpblane{2}{imp}{Improvement (I)}{0}{19.7}
  \hpblane{1}{imp}{Improvement (I)}{0}{19.7}
  \hpblane{-1}{iss}{Issue (X)}{0}{19.7}
  \hpblane{-2}{iss}{Issue (X)}{0}{19.7}
  \hpblane{-3}{task}{Task (T)}{0}{19.7}
  \hpblane{-4}{task}{Task (T)}{0}{19.7}
  % main line; history before v2.2 is collapsed into count badges
  \hpbmainline{0}{19}
  \hpbmainline*{19}{19.7}
  \hpborigin{0}{Product B}
  \hpbrelease{1}{v1.0}  \hpbcollapsed{1}{14}
  \hpbrelease{3}{v1.1}  \hpbcollapsed{3}{9}
  \hpbrelease{4}{v2.0}  \hpbcollapsed{4}{21}
  \hpbrelease{8}{v2.1}  \hpbcollapsed{8}{12}
  \hpbrelease{15}{v2.2}
  % long-lived open branches (outer lanes)
  \hpbbranch[label={T-14 EMC re-test}]{task}{1}{-4}{2}{19}{open}{}
  \hpbbranch[forkoff=0.6, label={L-5 Radio module}]{lab}{4}{4}{5}{19}{open}{}
  \hpbbranch[forkoff=0.4, label={T-20 Field logs}]{task}{4}{-3}{6}{19}{open}{}
  \hpbbranch[forkoff=0.7, label={X-88 Rare reboot}]{iss}{4}{-2}{7}{19}{open}{}
  % branches forked from v2.1 (nested rises: farther lane, smaller offset;
  % offsets at least 0.25 column apart and clear of the v2.1 tag)
  \hpbbranch[forkoff=0.55, label={L-7 Antenna tuning}]{lab}{8}{3}{9}{19}{open}{}
  \hpbbranch[forkoff=0.8, label={I-42 Fleet dashboard}]{imp}{8}{2}{10}{19}{open}{}
  \hpbbranch[forkoff=1.05, label={I-41 Quick setup}]{imp}{8}{1}{11}{13}{merge}{15}
  \hpbbranch[label={X-90 Clock drift}]{iss}{8}{-1}{12}{14}{merge}{15}
  % lanes reused after the merges into v2.2
  \hpbbranch[forkoff=0.6, label={I-45 Voice UI}]{imp}{15}{1}{16}{19}{open}{}
  \hpbbranch[label={X-95 Pairing}]{iss}{15}{-1}{17}{19}{open}{}
  % NOW marker (column 18 is left free so that labels clear the open heads)
  \hpbnow{19}{4}{-4}
  % callouts (boxes and leaders clear of every path and label)
  \hpbcallout{2.75,-2.3}{collapsed history:\\badge = merged\\branches hidden}{3,-0.62}
  \hpbcallout{2.35,3.2}{long-lived open\\branch, outer lane}{4.6,3.2}
  \hpbcallout{14.6,-4.85}{8 open branches at NOW on 8 lanes\\(= maximum overlap per type)}{19,-4.45}
\end{hpbgraph}}
\caption{Situation~E: many branches at once (illustrative). Older
releases are collapsed into count badges. Each type gets as many lanes as it
has overlapping branches, and freed lanes are reused: I-45 and X-95 take the
lanes of I-41 and X-90. Long-lived open work sits on the outer lanes; the
length of an open path counts the releases it has outlived, not days.}
\label{fig:scale}
\end{figure}


\paragraph{Situation.}
A mature product with years of history and many open branches. Draw
everything and the diagram becomes wallpaper. In a fast-rotating
portfolio this is the normal case.

\paragraph{In HPB.}
Two devices keep \cref{fig:scale} readable. Old releases are collapsed
into count badges, and lanes are reused: each type gets as many lanes as
it has branches open at the same time, stacked outward. Long-lived open
branches take the outer lanes, so newer forks reach the inner ones without
crossing them.

\paragraph{Decision.}
Which open work needs attention. Because time is counted in events, the
length of an open path shows how many releases it has outlived: T-14 has
survived four and goes on the next triage agenda. The lanes occupied at
NOW are the work in progress per type.

% =====================================================================
\subsection{F. The same items on a board and in HPB}
\label{sec:cond-board}

% F8 -- The same portfolio seen in a Kanban board and in HPB.
% Illustrative data. Requires hpb.sty v1.1. The board (a) is drawn with plain
% TikZ in explicit cm; the HPB view (b) uses the hpb macros and the document
% canvas. The panels are stacked so that both keep their natural size.
\begin{figure}[t]
\centering
\begin{tikzpicture}[
    kb/col/.style   = {draw=hpbDiscInk, line width=0.6pt, rounded corners=3pt,
                       fill=hpbGuide!30},
    kb/head/.style  = {font=\sffamily\bfseries\fontsize{7}{8.4}\selectfont,
                       text=hpbInk, anchor=north},
    kb/card/.style  = {draw=hpbInk, line width=0.8pt, rounded corners=2pt,
                       fill=white},
    kb/ref/.style   = {font=\sffamily\bfseries\fontsize{7}{8.4}\selectfont,
                       text=hpbInk, anchor=west, inner sep=0pt},
    kb/title/.style = {font=\sffamily\fontsize{6.5}{7.8}\selectfont,
                       text=hpbInk, anchor=west, inner sep=0pt},
    kb/panel/.style = {font=\sffamily\bfseries\fontsize{7.5}{9}\selectfont,
                       text=hpbInk, anchor=south west, inner sep=0pt},
  ]
  \node[kb/panel] at (0,1.95) {(a) Kanban board, Product A};
  % columns: x of left edge (cm), title
  \foreach \x/\t in {0/To do, 3.6/In progress, 7.2/Done}{
    \draw[kb/col] (\x,1.75) rectangle (\x+3.4,-1.0);
    \node[kb/head] at (\x+1.7,1.68) {\t};
    \draw[draw=hpbDiscInk, line width=0.6pt] (\x,1.3) -- (\x+3.4,1.3);
  }
  % cards: column left edge, top y (cm), type color, letter, ref, title
  \foreach \x/\y/\c/\l/\r/\t in {
      0/1.15/hpbIss/X/X-35/Drift,
      3.6/1.15/hpbLab/L/L-3/New sensor,
      3.6/0.35/hpbImp/I/I-14/Eco mode,
      7.2/1.15/hpbIss/X/X-31/Crash,
      7.2/0.35/hpbImp/I/I-12/Faster pairing,
      7.2/-0.45/hpbTask/T/T-7/Migrate logs}{
    \draw[kb/card] (\x+0.1,\y) rectangle (\x+3.3,\y-0.62);
    \node[hpb/event, fill=\c, minimum size=10pt, line width=1.1pt,
          font=\sffamily\bfseries\fontsize{6}{7}\selectfont]
      at (\x+0.38,\y-0.31) {\l};
    \node[kb/ref]   at (\x+0.65,\y-0.31) {\r};
    \node[kb/title] at (\x+1.35,\y-0.31) {\t};
  }
\end{tikzpicture}

\medskip
\hpbfit{%
\begin{hpbgraph}
  \node[font=\sffamily\bfseries\fontsize{7.5}{9}\selectfont, text=hpbInk,
        anchor=south west, inner sep=0pt] at (-2.75,3.0)
    {(b) HPB view of the same items};
  \hpblane{2}{lab}{Lab (L)}{0}{17.7}
  \hpblane{1}{imp}{Improvement (I)}{0}{17.7}
  \hpblane{-1}{iss}{Issue (X)}{0}{17.7}
  \hpblane{-2}{task}{Task (T)}{0}{17.7}
  \hpbmainline{0}{14}
  \hpbmainline*{14}{17.7}
  \hpborigin{0}{Product A}
  \hpbrelease{1}{v1.0}
  \hpbrelease{10}{v1.1}
  \hpbplanned{16}{v1.2}{Jun 2027}
  % history: two merges into v1.1, one discard
  \hpbbranch[forkoff=0.6, label={X-31 Crash}]{iss}{1}{-1}{2}{7}{merge}{10}
  \hpbbranch[forkoff=0.8, mergeoff=0.6, label={I-12 Faster pairing}]{imp}{1}{1}{3}{9}{merge}{10}
  \hpbbranch[forkoff=0.3, label={T-7 Migrate logs}]{task}{1}{-2}{4}{6}{discard}{}
  % open work at NOW
  \hpbbranch[forkoff=0.5, label={L-3 New sensor}]{lab}{1}{2}{5}{14}{open}{}
  \hpbstep{lab}{8}{2}{pass}
  \hpbbranch[forkoff=0.6, label={I-14 Eco mode}]{imp}{10}{1}{11}{14}{target}{16}
  \hpbbranch[label={X-35 Drift}]{iss}{10}{-1}{12}{14}{target}{16}
  \hpbnow{14}{2}{-2}
  % ordinal time axis with calendar ticks
  \hpbaxis{0}{17.7}{-2.75}
  \hpbtick{1}{Jan 2026}{-2.75}
  \hpbtick{10}{Jun}{-2.75}
  \hpbtick{14}{Nov}{-2.75}
\end{hpbgraph}}
\caption{Situation~F: the same six items on a Kanban board (a) and in
HPB (b) (illustrative). The board shows the workflow state of each item. HPB
adds what the board cannot show: the version each item started from, the
version that closed it, that T-7 was discarded rather than delivered, and
which items are committed to \texttt{v1.2} rather than merely open.}
\label{fig:same-data}
\end{figure}


\paragraph{Situation.}
Most teams already have a board, and they are right to. The first
question I am asked is: why do I need anything else?

\paragraph{In HPB.}
\Cref{fig:same-data} shows the same six items twice: on a three-column
Kanban board~\citep{anderson2010kanban} and as an HPB diagram. On the
board, X-31, I-12 and T-7 are all \emph{Done}. In HPB, X-31 and I-12 merge
into \texttt{v1.1} while T-7 was discarded; X-35 and I-14 are committed to
\texttt{v1.2}, and L-3 to nothing.

\paragraph{Decision.}
The board answers a team question, \emph{in which workflow state is each
item?}, and answers it well. HPB answers the product questions: what was
delivered and in which version, what is committed to the next gate, and
which hardware reference the work belongs to. Both views can be driven by
the same tracker items (\cref{subsec:method-coexistence}).

% =====================================================================
\subsection{G. Maintained versions and back-ports}
\label{sec:cond-backport}

% F11 -- Condition G: maintained release stream and back-port.
% Illustrative data. Requires hpb.sty v1.1. Canvas: document defaults.
\begin{figure}[t]
\centering
\hpbfit{%
\begin{hpbgraph}
  % ---- main stream of the line (y = 0) and its lanes
  \hpblane{1}{imp}{Improvement (I)}{0}{16.7}
  \hpblane{-1}{iss}{Issue (X)}{0}{16.7}
  \hpbmainline{0}{13}
  \hpbmainline*{13}{16.7}
  \hpborigin{0}{Gateway GW-2}
  \hpbrelease{1}{v1.4}
  \hpbrelease{6}{v2.0}
  \hpbrelease{12}{v2.0.1}
  % ---- maintenance stream 1.x (y = -2.2), forked at v1.4, own band
  \hpbseparator{-1.6}{-2.4}{16.7}
  \hpblane{-2.2}{main}{1.x stream\\(maintenance)}{0}{1.8}
  \hpblane{-3.2}{iss}{Issue (X)}{0}{16.7}
  \hpbstream[0.8]{1}{-2.2}{13}
  \hpbmainline*{13}{16.7}[-2.2]
  \hpbrelease{7}{v1.4.1}[-2.2]
  \hpbplanned{15}{v1.4.2}{Feb 2027}[-2.2]
  % ---- branches
  \hpbbranch[forkoff=0.6, label={I-70 Mesh networking}]{imp}{1}{1}{3}{5}{merge}{6}
  \hpbbranch[forkoff=0.3, mergeline=-2.2, label={X-298 Log rotation}]{iss}{1}{-3.2}{2}{4}{merge}{7}
  \hpbbranch[label={X-305 TLS parsing}]{iss}{6}{-1}{8}{11}{merge}{12}
  \hpbstep{iss}{9}{-1}{pass}
  \hpbbranch[line=-2.2, label={X-306 Port of X-305}]{iss}{7}{-3.2}{10}{13}{target}{15}
  % ---- propagation: X-306 carries the fix of X-305 to the 1.x stream (the
  %      connector reaches the start node from the upper left, clear of its label)
  \hpbpropagates{8,-1}{10,-3.2}{back-port}
  % ---- NOW and ordinal time axis
  \hpbnow{13}{1}{-3.2}
  \hpbaxis{0}{16.7}{-4.0}
  \hpbtick{1}{Mar 2026}{-4.0}
  \hpbtick{6}{Jul}{-4.0}
  \hpbtick{7}{Aug}{-4.0}
  \hpbtick{13}{Nov}{-4.0}
  \hpbtick{15}{Feb 2027}{-4.0}
\end{hpbgraph}}
\caption{Situation~G: a maintained stream and a back-port
(illustrative). The 1.x stream leaves \texttt{v1.4}; its patch
\texttt{v1.4.1} ships after \texttt{v2.0} but is not part of
\texttt{v2.0.1}. X-305 is fixed in \texttt{v2.0.1} and ported to 1.x as a
separate branch, X-306, linked by a propagation connector and targeted at
\texttt{v1.4.2}.}
\label{fig:backport}
\end{figure}


\paragraph{Situation.}
Some customers cannot move to the new major version, so the old one is
maintained in parallel. Every fix now raises a second question: does it
also have to go into the old version, and did anyone do it?

\paragraph{In HPB.}
In \cref{fig:backport}, a 1.x maintenance stream leaves \texttt{v1.4} in
its own band. Its patch \texttt{v1.4.1} ships after \texttt{v2.0}, but
precedence follows streams, not dates: no path connects them. X-305 is
fixed in \texttt{v2.0.1}; since 1.x is affected too, X-306 is opened on
the stream as a port, linked to X-305 by a propagation connector.

\paragraph{Decision.}
Which maintained versions contain a fix (follow the rises and the stream)
and which fixes have not been ported yet (merged issues known to affect
another stream, with no port and no recorded decision). The cost of
keeping 1.x alive is the work that forks from it.

% =====================================================================
\subsection{H. From field return to corrective action}
\label{sec:cond-rma}

% F12 -- Condition H: field return to corrective action.
% Illustrative data. Requires hpb.sty v1.1. Canvas: document defaults.
% No improvement lane in this diagram, so the lab lane is the first lane
% above the main line (lanes are stacked outward without gaps).
\begin{figure}[t]
\centering
\hpbfit{%
\begin{hpbgraph}
  \hpblane{1}{lab}{Lab (L)}{0}{17.7}
  \hpblane{-1}{iss}{Issue (X)}{0}{17.7}
  \hpblane{-2}{iss}{Issue (X)}{0}{17.7}
  \hpbmainline{0}{13}
  \hpbmainline*{13}{17.7}
  \hpborigin{0}{DMS-4000\\Rev B}
  \hpbrelease{1}{v2.0}
  \hpbprocess[]{8}{PC-7}
  \hpbplanned[interchangeable]{15}{Rev B.1}{Mar 2027}
  % field return closed without a design change: contained (rationale class
  % "resolved without a product baseline")
  \hpbbranch[forkoff=0.6, label={X-152 Housing damage}]{iss}{1}{-1}{3}{4}{discard}{}
  \hpblabel[right]{4}{-1}{\itshape contained:\\\itshape stock reworked}
  % field return with failure analysis, closed by a process baseline, then
  % checked for effectiveness after closure
  \hpbbranch[forkoff=0.3, label={X-150 Relay failure (RMA)}]{iss}{1}{-2}{2}{7}{merge}{8}
  \hpbstep{iss}{5}{-2}{done}
  \hpbstep{iss}{6}{-2}{pass}
  \hpbpoststep{iss}{7.5}{10}{-2}{pass}
  \hpblabel{10}{-2}{\itshape effectiveness verified}
  % supplier end-of-life notice: qualification of an alternative relay,
  % targeted at an interchangeable hardware baseline
  \hpbbranch[forkoff=0.6, label={L-9 PCN-221 relay EOL}]{lab}{8}{1}{9}{13}{target}{15}
  \hpbstep{lab}{11}{1}{pass}
  \hpbnow{13}{1}{-2}
  % ordinal time axis
  \hpbaxis{0}{17.7}{-3.0}
  \hpbtick{1}{Apr 2026}{-3.0}
  \hpbtick{8}{Jul}{-3.0}
  \hpbtick{13}{Nov}{-3.0}
  \hpbtick{15}{Mar 2027}{-3.0}
\end{hpbgraph}}
\caption{Situation~H: from field return to corrective action
(illustrative). X-150 is solved by changing how the product is built, so it
merges into the process baseline PC-7 (square), and a step after closure
records that the action worked. X-152 is closed as contained, without a
product baseline. The supplier notice PCN-221 opens L-9, targeted at the
interchangeable hardware baseline \texttt{Rev~B.1}.}
\label{fig:rma}
\end{figure}


\paragraph{Situation.}
Units come back from the field. Some failures are fixed by changing how
the product is built, not what it is; some are not product problems at
all. Quality wants to know what was done, on which evidence, and whether
it worked.

\paragraph{In HPB.}
In \cref{fig:rma}, X-150 \emph{Relay failure} carries the failure analysis
and the verification on a pilot lot as steps, and merges into the process
baseline PC-7 (a square). A dotted continuation records that the action
was effective. X-152 closes as \emph{resolved without a product baseline}
(transport damage, stock reworked). A supplier end-of-life notice opens
the lab branch L-9, targeted at an interchangeable hardware baseline.

\paragraph{Decision.}
The corrective-action loop is recorded without inventing product
versions. Containment is its own closure class and does not inflate the
abandoned-work ratio; effectiveness is a step after closure.

% =====================================================================
\subsection{I. Shared firmware and a joint launch}
\label{sec:cond-shared}

% F13 -- Condition I: firmware shared by two hardware references, with a
% dependency and a coordinated launch.  Illustrative data.  Requires hpb.sty
% v1.1.  Canvas: document defaults.  Three lines, each in its own band:
% hardware Rev B (top), the shared firmware line (middle), hardware Rev A.
\begin{figure}[t]
\centering
\hpbfit{%
\begin{hpbgraph}
  % ---- hardware line DMS-4000 Rev B (y = 2.7)
  \hpbmainline{0}{10}[2.7]
  \hpbmainline*{10}{17.7}[2.7]
  \hpborigin{0}{DMS-4000\\Rev B}[2.7]
  \hpbrelease{2}{B.1}[2.7]
  \hpbplanned{13}{B.2}{Apr 2027}[2.7]
  \hpbseparator{1.85}{-2.4}{17.7}
  % ---- shared firmware line (y = 0) with its lanes
  \hpblane{1}{imp}{Improvement (I)}{0}{17.7}
  \hpblane{-1}{iss}{Issue (X)}{0}{17.7}
  \hpbmainline{0}{10}
  \hpbmainline*{10}{17.7}
  \hpborigin{0}{DMS firmware}
  \hpbrelease{3}{v2.0}
  \begin{pgfonlayer}{hpbfront}% applies-to set, left of the tag (right side
    % is reserved for upward forks)
    \node[font=\sffamily\itshape\fontsize{6.5}{7.5}\selectfont, text=hpbMuted,
          anchor=east, inner sep=1pt] at ([xshift=-1.5pt]hpb@t.west)
      {applies to A, B};
  \end{pgfonlayer}
  \hpbplanned[applies to A, B]{14}{v2.1}{Apr 2027}
  \hpbseparator{-1.85}{-2.4}{17.7}
  % ---- hardware line DMS-4000 Rev A (y = -2.7)
  \hpbmainline{0}{10}[-2.7]
  \hpbmainline*{10}{17.7}[-2.7]
  \hpborigin{0}{DMS-4000\\Rev A}[-2.7]
  \hpbrelease{1}{A.3}[-2.7]
  \hpbmilestone{7}{end of sale}[-2.7]
  % ---- branches of the firmware line
  \hpbbranch[forkoff=0.5, label={X-140 Relay chatter (Rev A + v2.0)}]{iss}{3}{-1}{4}{10}{target}{14}
  \hpbstep{iss}{8}{-1}{pass}
  \hpbbranch[forkoff=0.6, label={I-77 Low-power mode}]{imp}{3}{1}{5}{10}{target}{14}
  \hpbstep{imp}{6}{1}{done}
  % ---- I-77 depends on the hardware baseline B.2; launch bundle
  \hpbdepends{10,1}{13,2.7}{depends on}
  \hpbbundle{12.1,-0.55}{17.4,3.85}{launch bundle K-1 (all or nothing)}
  \hpbnow{10}{2.7}{-2.7}
  % ---- ordinal time axis
  \hpbaxis{0}{17.7}{-3.75}
  \hpbtick{1}{Feb 2026}{-3.75}
  \hpbtick{3}{May}{-3.75}
  \hpbtick{7}{Sep}{-3.75}
  \hpbtick{10}{Nov}{-3.75}
  \hpbtick{13}{Apr 2027}{-3.75}
\end{hpbgraph}}
\caption{Situation~I: one firmware shared by two hardware references
(illustrative). Each firmware version declares the hardware it applies to;
X-140 is drawn on the firmware line, with its observed configuration in the
label. I-77 depends on the hardware baseline B.2, so \texttt{v2.1} and B.2
form the launch bundle K-1: published together or not at all
(\cref{rule:gate}).}
\label{fig:shared}
\end{figure}


\paragraph{Situation.}
One firmware build runs on several hardware references, and a new feature
needs both a firmware release and a hardware change. A report says
``relay chatter'', but the real question is: on which combinations?

\paragraph{In HPB.}
In \cref{fig:shared}, the firmware is a line of its own between the
hardware lines Rev~A and Rev~B, and each firmware version declares the
hardware it applies to. X-140 is drawn on the firmware line with its
observed configuration (Rev~A, \texttt{v2.0}). I-77 depends on the
hardware baseline B.2, so \texttt{v2.1} and B.2 form the launch bundle
K-1.

\paragraph{Decision.}
Who is affected (the applies-to sets and the observed configuration) and
whether the launch can go ahead: by \cref{rule:gate}, K-1 ships only when
everything targeted at either baseline is verified. A slip of B.2 visibly
delays \texttt{v2.1}, unless I-77 is moved out.
